\clearpage\phantomsection
\addcontentsline{toc}{chapter}{DAFTAR PUSTAKA}
\begin{thebibliography}{99}
\bibitem{spesifikasi}
A. E. R. Xavier dan N. G. Subrata,
``Tugas 2 IF4020-Kriptografi Semester 1 Tahun 2026/2027:
Membuat Block Cipher Baru,'' spesifikasi tugas, 2026.
Salinan lokal: \path{Tugas2-IF4020-2026.md}.

\bibitem{aes}
National Institute of Standards and Technology,
``Advanced Encryption Standard (AES),'' FIPS 197-upd1,
9 Mei 2023. [Daring]. Tersedia:
\url{https://doi.org/10.6028/NIST.FIPS.197-upd1}.

\bibitem{mode}
M. Dworkin, ``Recommendation for Block Cipher Modes of Operation:
Methods and Techniques,'' NIST SP 800-38A, Desember 2001.
[Daring]. Tersedia: \url{https://doi.org/10.6028/NIST.SP.800-38A}.

\bibitem{padding}
R. Housley, ``Cryptographic Message Syntax (CMS),'' RFC 5652,
bagian 6.3 dan 9, September 2009. [Daring]. Tersedia:
\url{https://www.rfc-editor.org/rfc/rfc5652}.

\bibitem{shannon}
C. E. Shannon, ``Communication Theory of Secrecy Systems,''
\textit{Bell System Technical Journal}, vol. 28, no. 4,
hlm. 656--715, Oktober 1949. [Daring]. Tersedia:
\url{https://doi.org/10.1002/j.1538-7305.1949.tb00928.x}.

\bibitem{heys}
H. M. Heys, ``A Tutorial on Linear and Differential Cryptanalysis,''
laporan teknis CORR 2001-17, Centre for Applied Cryptographic Research,
University of Waterloo, Maret 2001. [Daring]. Tersedia:
\url{https://www.engr.mun.ca/~howard/PAPERS/ldc_tutorial.pdf}.

\bibitem{cryptosys}
D. Ireland, ``Padding schemes for block ciphers,''
\textit{CryptoSys PKI Pro Manual}, DI Management Services,
2026. [Daring]. Tersedia:
\url{https://www.cryptosys.net/pki/manpki/pki_paddingschemes.html}.
\end{thebibliography}

% Tambahkan sumber yang benar-benar dibaca/digunakan; kutip dengan \cite{...}.
